% SOURCE_STATEMENT: "1. When M is a 3-manifold are there some 2-dimensional objects that are related to Spec(∆) in a similar way? 2. Does there exist some analogue for the \"area spectrum\" of minimal surfaces in M 3 (with some suitable geometric assumptions imposed on M 3)? 3. Compute the number of minimal surfaces N (n1, n 2, g ) in T 4 = T 2 × T 2 endowed with the flat metric which represent the homology class n1[T1] + n2[T2], where Ti are the two 2-tori factors of T 4."
% Preparation and final self-audit: CONTRIBUTING.md prompts C and D; source/audit artifacts retained in the local physics_statements_2026_10_11 batch.
% !TeX program = lualatex
\documentclass[11pt]{article}
\usepackage[margin=25mm]{geometry}
\usepackage{fontspec}
\setmainfont{Latin Modern Roman}
\usepackage{amsmath,amssymb,mathtools}
\usepackage{unicode-math}
\setmathfont{Latin Modern Math}
\usepackage{xurl,hyperref}
\hypersetup{colorlinks=true,urlcolor=blue}
\setcounter{secnumdepth}{0}
\setlength{\parindent}{0pt}
\setlength{\parskip}{5pt}
\setlength{\emergencystretch}{3em}
\newcommand{\sref}[2]{\hyperlink{source:#1:#2}{[S#2]}}
\newcommand{\cataloguescope}[1]{\paragraph{Recorded target.}#1}
\newcommand{\R}{\mathbb R}
\newcommand{\C}{\mathbb C}
\newcommand{\N}{\mathbb N}
\newcommand{\Z}{\mathbb Z}
\newcommand{\Q}{\mathbb Q}
\newcommand{\F}{\mathbb F}
\newcommand{\eps}{\varepsilon}
\begin{document}
\section{Audible primitive two-torus areas on a flat three-torus}\label{statement}

\subsection{Definitions and mathematical statement}

\paragraph{Shared notation.}$\mathbb N=\{1,2,\ldots\}$; the standard number fields have their usual meanings. All additional conventions are stated locally.


\paragraph{Definitions and notation.}$\operatorname{Spec}(\Delta)$ is the spectrum of the Laplace operator for the chosen Riemannian metric and boundary condition. Minimal surfaces are critical points of area. An area spectrum requires specifying which minimal surfaces are included and how multiplicities or families are counted. In $T^4=T^2\times T^2$, $[T_1],[T_2]$ are the homology classes of the two torus factors; $g$ denotes the surface genus.

\paragraph{Catalogue wording retained for traceability.}
For three-manifolds, identify two-dimensional geometric objects related to $\operatorname{Spec}(\Delta)$ in a manner analogous to the geodesic-length relationship. Determine whether a suitable analogue of the minimal-surface area spectrum exists under appropriate geometric hypotheses. Compute the number $N(n_1,n_2,g)$ of genus-$g$ minimal surfaces in the flat $T^4=T^2\times T^2$ representing $n_1[T_1]+n_2[T_2]$.
\paragraph{Scope and source qualification.}The source leaves the counting convention and the relevant area-spectrum objects open; translations and continuous families cannot silently be discarded. The title’s primitive two-torus example does not replace the full three-part target. The extracted background remarks spill into separate subsequent Jardim Dirac-spectrum questions and the Hausel programme; those are not silently reassigned to this canonical three-part Cherkis excerpt. The catalogue title/status was paired with a machine-generated research attempt; that attempt is not peer-reviewed evidence of settlement. The statement covers the original source questions rather than only the title’s subcase.
\subsection{Short English statement}
Find surface-area analogues of spectral/length relations and count minimal surfaces in prescribed flat four-torus classes.
\subsection{Sources}
\begin{itemize}
\item[\hypertarget{source:statement:1}{S1}] ULAM.AI and the UnsolvedMath Contributors, UnsolvedMath, catalogue entry. Catalogue curation under CC BY 4.0; definitions and English exposition added locally.
\url{https://huggingface.co/datasets/ulamai/UnsolvedMath}.
\item[\hypertarget{source:statement:2}{S2}] American Institute of Mathematics, $L^2$ harmonic forms in geometry and string theory; notes by Anda Degeratu and Mark Haskins, March 2004, 18 March 2004, Cherkis area-spectrum questions.
\url{https://aimath.org/WWN/l2harmonic/l2harmonic.pdf}.
\end{itemize}
\label{end:statement}
\end{document}
